% =============================================================================
% --- secciones/2-Diagnostico.tex ---
% =============================================================================
\section{Diagnóstico}

\subsection{La Brecha entre el Urbanismo Teórico y la Realidad Social}

El diagnóstico de la movilidad en las metrópolis mexicanas revela una desconexión profunda 
entre los modelos de planeación indicativa y la dinámica físico-social del territorio. 
Mientras la teoría urbana contemporánea promueve la consolidación de ciudades compactas y 
resilientes, la praxis del desarrollo habitacional y la infraestructura de transporte ha 
fomentado un modelo de dispersión que compromete la viabilidad sistémica de la región 
\cite{NOM004, ConstitucionMex}.

\subsection{La Ciudad Expandida: El Fracaso de la Compacidad}

Uno de los mayores desafíos para la sustentabilidad es la expansión anárquica de las manchas 
urbanas. Datos históricos de SEDESOL indican que las ciudades en México han crecido físicamente 
a tasas desproporcionadas en comparación con su demografía: mientras la superficie urbana se 
ha incrementado \textbf{6 veces}, la población solo ha crecido \textbf{1.9 veces} 
\cite{ITDP_Espacio, IMCO2019Doc}. 

Este fenómeno de baja densidad no es un proceso neutral; representa un modelo disperso que 
favorece la ineficiencia energética y profundiza la segregación social. Al agotarse 
el suelo central asequible, las poblaciones de menores ingresos son expulsadas hacia 
periferias distantes que carecen de servicios de transporte masivo de calidad, 
forzándolas a depender de redes capilares informales o de baja capacidad 
\cite{ramirez2025analisis, RedalycPeaton}.

\subsection{Motorización Desenfrenada y Desequilibrio Estadístico}

La crisis de movilidad se ve exacerbada por una tendencia de motorización que anula 
cualquier ganancia en eficiencia vial. De acuerdo con el Instituto Mexicano para la 
Competitividad (IMCO), el parque vehicular en las zonas metropolitanas crece a una 
tasa promedio anual del \textbf{5.3\%}, un ritmo \textbf{3.5 veces más rápido} que el 
crecimiento poblacional (1.5\%) \cite{Litman2021, batista2015}. 

Este crecimiento no se traduce en una mayor capacidad de desplazamiento para la colectividad, 
sino en una ocupación desmedida del espacio público. En la ZMVM, se estima que el 
\textbf{68.1\% de los viajes en automóvil particular} transportan a un solo pasajero, lo 
que genera externalidades negativas (congestión, siniestralidad y contaminación) que 
representan un costo equivalente al \textbf{4\% del PIB nacional} 
\cite{IMCO2019Doc, RedalycPeaton}. La ineficiencia del automóvil 
es el síntoma de un sistema que prioriza la velocidad del objeto sobre la accesibilidad 
de la persona.

\subsection{El Hallazgo VFT: Segregación en el Cuarto Contorno}

Desde la perspectiva del \textbf{Modelo VFT}, la brecha entre la teoría y la realidad 
se manifiesta con mayor crudeza en la periferia profunda. El análisis del 
\textbf{Nivel de Cobertura ($C$)} demuestra que el diseño radial del transporte 
masivo ha agotado su capacidad operativa al llegar al tercer y cuarto contorno 
de la ciudad \cite{SEDATU2023, Garibaldi2024VFT}. 

Los resultados del diagnóstico espacial evidencian que demarcaciones periféricas 
presentan niveles de cobertura marginales inferiores al \textbf{12\%}, siendo el caso 
de Milpa Alta (11.75\%) el más crítico \cite{Garibaldi2024VFT}. 
Esta realidad social invalida el principio de accesibilidad universal consagrado en el 
Artículo 4° Constitucional y reduce a los habitantes de estas zonas a ciudadanos de 
"segunda o tercera clase", segregados del derecho a la ciudad por una barrera de distancia 
y tiempo \cite{rendon2009ciudadanos}. Esta exclusión justifica la necesidad de una 
reconfiguración topológica que rompa la inercia radial para transitar hacia una red de 
interceptación transversal.
